\rok{Јануар 2024.}{H}

Други практични тест из Алгоритама и структура података

Други практични тест одржан 12.01.2024.

\zadatak{1}{Конверзија матрице у листу суседства}
Написати алгоритам који ће вршити директну конверзију матрице суседства у листу суседства. Алгоритам као улазни параметар треба да прихвати матрицу, а треба да резултује креираним динамичким низом чворова.

\textbf{Додатне напомене за решавање задатка:}
\begin{itemize}
	\item Написати алгоритам.
	\item У Main методи креирати произвољан граф. Обезбедити начин за тестирање и проверу нове функционалности, односно позив \texttt{convertMatrixToList} методе.
	\item У template-у је закоментарисана метода \texttt{convertMatrixToList} која треба да садржи имплементацију конверзије листе у матрицу суседства.
\end{itemize}



\zadatak{2}{Исписивање свих непарних вредности листова}
У оквиру стандардне структуре класе \texttt{AVL} написати имплементацију методе која ће омогућити проналажење и исписивање чворова на позицији листова са непарном вредношћу.

\textbf{Додатне напомене за решавање задатка:}
\begin{itemize}
	\item У оквиру класе \texttt{AVL} креирати јавну методу са потписом:
	\begin{Verbatim}[fontsize=\footnotesize, numbers=left, xleftmargin=10mm]
public void printOddLeaves()
	\end{Verbatim}
	\item Из претходно креиране јавне методе треба да се позива приватна метода са потписом:
	\begin{Verbatim}[fontsize=\footnotesize, numbers=left, xleftmargin=10mm]
private Node printOddLeavesAVL(Node curr)
	\end{Verbatim}
	\item Приватна метода треба да садржи имплементацију алгоритма којим ће се омогућити пролазак кроз стабло од корена ка листовима и њихово исписивање уколико им је непарна вредност.
	\item Листом се класификује сваки чвор који нема ниједно дете.
	\item Може се креирати помоћна променљива на нивоу класе \texttt{AVL} која ће помоћи у пребројавању.
	\item Алгоритам \textbf{ОБАВЕЗНО} писати у оквиру класе \texttt{AVL}.
\end{itemize}

\textbf{Примери:}

\begin{center}
\begin{minipage}{\textwidth}
\centering
\begin{forest}
for tree={
	circle,
	draw,
	thick,
	fill=gray!20,
	minimum size=8mm,
	inner sep=1pt,
	s sep=8mm,
	l sep=12mm,
	edge={-, thick},
	font=\small
}
[8
	[4
		[2
			[, phantom]
			[3]
		]
		[5]
	]
	[9
		[, phantom]
		[11]
	]
]
\end{forest}

\vspace{0.3cm}
\textbf{Слика 1.} Пример AVL-а.
\end{minipage}
\end{center}

\begin{itemize}
	\item \textbf{Улаз 1:} AVL са слике 1
	\item \textbf{Излаз 1:} 3, 5, 11
\end{itemize}

\begin{center}
\begin{minipage}{\textwidth}
\centering
\begin{forest}
for tree={
	circle,
	draw,
	thick,
	fill=gray!20,
	minimum size=8mm,
	inner sep=1pt,
	s sep=8mm,
	l sep=12mm,
	edge={-, thick},
	font=\small
}
[10
	[4
		[2
			[3]
			[, phantom]
		]
		[6
			[5]
			[9]
		]
	]
	[20
		[12
			[, phantom]
			[14]
		]
		[29
			[27]
			[, phantom]
		]
	]
]
\end{forest}

\vspace{0.3cm}
\textbf{Слика 2.} Пример AVL-а.
\end{minipage}
\end{center}

\begin{itemize}
	\item \textbf{Улаз 2:} AVL са слике 2
	\item \textbf{Излаз 2:} 3, 5, 9, 27
\end{itemize}



\zadatak{3}{Провера могућности мутације подсеквенце}
Написати алгоритам који проверава да ли је могуће мутирати секвенцу. Алгоритам као параметре добија секвенцу за коју се проверава да ли је могућа мутација, подсеквенцу коју треба мутирати и вредност мутације. Секвенцу је могуће мутирати уколико је естимирана сличност секвенце пре мутације са секвенцом након мутације бар 0,5.
Алгоритам написати уз претпоставку да се тражена подсеквенца коју је потребно мутирати појављује највише једанпут у секвенци.

\textbf{Додатне напомене за решавање задатка:}
\begin{itemize}
	\item Написати алгоритам.
	\item У Main методи обезбедити тестирање, позивом методе са параметрима који су дати у примерима.
	\item Користити \texttt{search} методу класе \texttt{KarpRabinAlgorithm} и класе \texttt{MinHash} и \texttt{Jaccard}, које су дате у шаблону.
	\item Имплементацију писати у оквиру закоментарисане методе:
	\begin{Verbatim}[fontsize=\footnotesize, numbers=left, xleftmargin=10mm]
public static bool checkIfMutationIsPossible(string str, string mutatingSubsequence, string mutationSubsequence)
	\end{Verbatim}
	\item Приликом креирања \texttt{MinHash} објекта, за вредности параметара \texttt{universeSize} и \texttt{numHashFunctions} прослеђивати вредности 12 и 4.
	\item Параметар \texttt{q} функције \texttt{search} поставити на \texttt{mutatingSubsequence.Length+1}.
	\item Дозвољено је изменити повратни тип методе \texttt{search} (нпр. уместо \texttt{void} имати \texttt{List<int>}).
\end{itemize}

\textbf{Напомене које могу бити од помоћи приликом решавања задатка:}
\begin{itemize}
	\item За конверзију string-а у низ карактера користити \texttt{ToCharArray()} по принципу: \texttt{str.ToCharArray()}.
	\item За конверзију низа карактера у string користити \texttt{new string(char[] arrayOfChars)} по принципу: \texttt{new string(charArray)}.
	\item Карактере није потребно експлицитно конвертовати у целобројни тип приликом убацивања у \texttt{List<uint>} структуру.
\end{itemize}

\textbf{Пример:}
\begin{itemize}
	\item \textbf{Улаз 1:} \texttt{"ACGTWMAB", "CGT", "GTB"}
	\item \textbf{Излаз 1:} \texttt{True}
	\item \textbf{Образложење:} Мутацијом се добија секвенца \texttt{AGTBWMAB}, сличност секвенци \texttt{ACGTWMAB} и \texttt{AGTBWMAB} је преко 0.5.
	\item \textbf{Улаз 2:} \texttt{"AMCGT", "CGT", "HBK"}
	\item \textbf{Излаз 2:} \texttt{False}
	\item \textbf{Образложење:} Мутацијом се добија секвенца \texttt{AMHBK}, сличност секвенци \texttt{AMCGT} и \texttt{AMHBK} није преко 0.5.
\end{itemize}



\sablon{Шаблон другог теста јануар 2024. група H}{2023/Januar/GrupaH/AISP2023_Test2_GrupaH.linq}
